\documentclass[11pt]{article}
\usepackage[margin=1in]{geometry}
\usepackage{amsmath,amssymb}
\usepackage[hidelinks]{hyperref}

\begin{document}

\section*{Human Review: Positive multilinear forms on non-atomic \texorpdfstring{$L_p$}{Lp}}

\noindent Original automatic proof: \texttt{solution\_packet.pdf}

\noindent Checked date: 13 June 2026

\noindent Status: Manually checked.

\noindent Verdict: The proof appears correct.

\section*{Human Comments}

The proposed solution appears to give a correct negative answer to Question 3
from the source paper, under the natural interpretation that the spaces are in
the finite-exponent range \(1<p_i<\infty\). The counterexample is the positive
\(n\)-linear form
\[
  A(f_1,\ldots,f_n)=\int_0^1 f_1(t)\cdots f_n(t)\,dt
\]
on
\[
  L_{p_1}[0,1]\times\cdots\times L_{p_n}[0,1],
  \qquad
  \sum_{i=1}^n \frac1{p_i}<1.
\]
This form is bounded and positive, but it is not weakly sequentially
continuous.

\section*{Boundedness and Positivity}

The boundedness is the standard Holder argument. Since
\[
  \sum_{i=1}^n \frac1{p_i}<1,
\]
one can choose \(r\in(1,\infty]\) such that
\[
  \sum_{i=1}^n \frac1{p_i}+\frac1r=1.
\]
Applying Holder's inequality to \(f_1\cdots f_n\cdot 1\) gives
\[
  |A(f_1,\ldots,f_n)|
  \leq
  \|f_1\|_{p_1}\cdots \|f_n\|_{p_n}\|1\|_r.
\]
Thus \(A\) is a continuous \(n\)-linear form. Positivity is immediate: if
each \(f_i\geq0\), then the integrand is non-negative almost everywhere.

\section*{Failure of Weak Sequential Continuity}

Let \((r_k)\) be the Rademacher functions. For every finite \(p\), the
sequence \((r_k)\) is weakly null in \(L_p[0,1]\). Hence \(r_k\to0\) weakly in
each relevant finite \(L_{p_i}\)-space.

However,
\[
  A(r_k,r_k,1,\ldots,1)
  =
  \int_0^1 r_k(t)^2\,dt
  =
  1
\]
for every \(k\), whereas
\[
  A(0,0,1,\ldots,1)=0.
\]
Therefore \(A\) is not weakly sequentially continuous.

\section*{Interpretation}

The argument is short, but it seems to be a genuine obstruction to the stated
question. The main point to check is not the calculation itself, but whether
the authors intended the question exactly in this finite-exponent form, or
whether some additional hypothesis was meant to exclude this Rademacher-product
example.

\section*{Review Outcome}

Archive this packet as an apparently correct full negative answer, pending
author feedback. I recommend contacting the authors to ask whether the
argument is known to them, whether it matches the intended formulation of the
question, and how they would assess the value of the example.

\end{document}
